% ===== tr-08 · b15 · tr08f：第 5 章 p41--42 + p43 顶部正文（原文页 41--43）=====
% [原文 p.41]
在我们的应用中，psh 函数 $\phi$ 来自 $X_t$ 内部环面坐标卡上 Calabi--Yau 度量 $\omega_{CY,t}$ 的局部势 $\phi_{CY,J,t}$。由于 $\omega_{CY,t}$ 相对于 Fubini--Study 参考度量具有一致有界的势（参看定理 5.4.6），容易安排环面坐标卡上的局部势 $\phi_{CY,J,t}$ 一致有界，从而凸函数 $\bar\phi_{CY,J,t}$ 也一致有界。利用 Lipschitz 界，由 Arzela--Ascoli 定理，我们可以当 $t\to 0$ 时抽出一族子序列极限。按构造，在 $Sk(X)$ 各 $n$ 维面的内部有 $\bar\phi_{CY,J,t}\to\bar\phi_{J,0}$（$C^0_{loc}$ 意义）。

% [Tutor Note: 本页起 \bar\phi_{CY,J,t} 的横线只盖住 $\phi$（共轭/平均记号），不盖下标 $CY,J,t$；$\bar\phi_{J,0}$ 的下标只有 $J,0$（没有 $CY$），二者不是同一个对象：前者是逐 $t$ 的凸平均，后者是极限。Arzela--Ascoli 这里用的是「一致有界 + Lipschitz 界」给出紧性，故只能在子序列上取极限。]

\subsubsection{势的 $C^0$-收敛与延拓问题}
我们希望证明：在略微收缩的 $U_{J,t}$ 上，$\|\phi_{CY,J,t}-\bar\phi_{J,0}\circ\mathrm{Log}_{\mathcal X}\|_{C^0}$ 沿该子序列收敛到零。我们知道沿子序列 $\|\bar\phi_{CY,t}-\bar\phi_{J,0}\|_{C^0}\to 0$，并且由注 5.17 知道 $|\phi_{CY,J,t}-\bar\phi_{CY,J,t}\circ\mathrm{Log}_{\mathcal X}|$ 除在一个小测度子集上之外都很小。然而，要去掉这个小测度问题相当微妙，需要整体的论证。

% [Tutor Note: 上一段 $\|\bar\phi_{CY,t}-\bar\phi_{J,0}\|_{C^0}$ 中的第一项是 $\bar\phi_{CY,t}$（无 $J$），与前面带 $J$ 的 $\bar\phi_{CY,J,t}$ 不同------这是原文的记号松弛，按页图照抄，读者应理解为在面 $\Delta_J$ 上的平均量。]

策略如下：
\begin{itemize}
\item （「延拓问题」）在 $X_t$ 上找一个整体 Kähler 度量 $\omega_{\psi,t}$，使其在 $U_{J,t}$ 上的局部势与 $\bar\phi_{J,0}\circ\mathrm{Log}_{\mathcal X}$ 相差 $C^0$ 小误差，且它相对于某个固定 Fubini--Study 度量的势与 $t$ 无关地有界；
\item （势稳定性估计）然后可以考察 Calabi--Yau 度量 $\omega_{CY,t}$ 相对于 $\omega_{\psi,t}$ 的势 $\phi_{CY,rel}$。$U_{J,t}$ 上 $\phi_{CY,rel}$ 的小上界来自 $\phi_{CY,J,t}-\bar\phi_{CY,J,t}\circ\mathrm{Log}_{\mathcal X}\le\frac{C}{|\log|t||^{1/2}}$。我们还知道 $\phi_{CY,rel}$ 的小下界除在一个测度极小的集合上之外成立，于是应用定理 5.4.7 便得到 $\inf\phi_{CY,rel}$ 的小下界。我们强调，Kähler 度量的整体正性对这一论证至关重要。
\end{itemize}

最终结论是：在 $U_{J,t}$ 略微收缩后的版本上，$\phi_{CY,rel}$ 是 $C^0$-小的。这等价于 $\|\phi_{CY,J,t}-\bar\phi_{J,0}\circ\mathrm{Log}_{\mathcal X}\|_{C^0}$ 的小性，而这正是我们的目标。

% [Tutor Note: $\phi_{CY,rel}$ 是「相对势」，即 $\omega_{CY,t}$ 相对于 $\omega_{\psi,t}$ 的势（$\omega_{CY,t}=\omega_{\psi,t}+dd^c\phi_{CY,rel}$ 型），不是新的度量。这里逻辑链是：延拓给出正则的整体 Kähler 度量 $\omega_{\psi,t}$，再用定理 4.7（$L^1$-稳定性）把相对势夹到小，从而得到 $C^0$ 收敛。「整体正性必不可少」------这是与 (11) 型 Skoda 估计衔接的前提。]

延拓问题关乎把局部势拼接成整体 Kähler 势，困难在于在非一般区域实现 psh 性质。其核心问题在整个一般情形下尚未令人满意地解决，即

\begin{question}{5.8}
我们能否足够显式地刻画 $Sk(X)$ 上那些可以看作 $X_t$ 上 Kähler 势之极限的凸势类？
\end{question}

% [Tutor Note: 问题编号原文为 Question 8，本章内套用 $\{5.8\}$。注意它问的不是「凸函数是否存在」而是「能否显式刻画这一类」------这正是对实 MA 方程作直接表述（以及后文处理仿射结构/奇异集）的前置条件。]

对 Fermat 族，延拓问题以一种特设（ad hoc）的方式得到解决，这是 [52] 中最技术性的部分。\footnote{严格说来，文献 [52] 并未使用对偶复形与本质骨架的语言，而是通过由热带几何控制的显式坐标卡来进行。}回忆 Fermat 族：它嵌入到环境射影空间 $\C\mathrm{P}^{n+1}$ 中。

% [原文 p.42]
我们的策略是把延拓 $\omega_{\psi,t}$ 构造为 $\C\mathrm{P}^{n+1}$ 上的环面 Kähler 度量，再限制到 $X_t$，这就保证了整体正性。要确保 $\omega_{\psi,t}$ 与局部凸函数相差 $C^0$ 小误差是一件精细的事：它涉及显式的热带超曲面组合学，利用 Fermat 族大量的离散对称性，并使用一个双重 Legendre 变换构造 [52]。

% [Tutor Note: 脚注 5 是理解策略 II 的关键：文献 [52] 的原始论证其实不用对偶复形/本质骨架，而用热带几何控制的显式坐标卡，因此本节的「$Sk(X)$／双复形」语言是对 [52] 的再翻译。$\C\mathrm{P}^{n+1}$ 在原文中排作 ℂ（双线体）＋正体 P 的上标，不是两个字母都双线体。]

\begin{remark}{5.18}
在 $\C\mathrm{P}^{n+1}$ 上考虑环面 Kähler 度量的动机如下：环面性质是把一般 Kähler 势化为凸函数的自然途径；而延拓到环境空间的想法则以如下结果为据：
\end{remark}

\begin{proposition}{5.6.8}[19, Thm. B]
设 $(X,\omega)$ 是射影流形，其 Kähler 形式代表一个整类，$Y$ 是 $X$ 的光滑子簇。则任何 $\phi\in PSH(Y,\omega|_Y)$ 都延拓为 $\phi\in PSH(X,\omega)$。
\end{proposition}

% [Tutor Note: 命题 6.8 的 $PSH(Y,\omega|_Y)$ 中 $\omega|_Y$ 是 $\omega$ 到子簇 $Y$ 的限制（竖线下标 $Y$），三个字母 $PSH$ 为数学斜体（页图如此），不要排成正体。该命题是「环境空间延拓」的一般依据；在 Fermat 情形取 $(X,Y)=(\C\mathrm{P}^{n+1},X_t)$。]

\subsubsection{实 MA 度量}
为了补全这一循环，我们需要

\begin{lemma}{5.6.9}
极限凸势 $\bar\phi_{J,0}$ 在 $\Delta_J$ 的内部解实 MA 方程 (5.17)。
\end{lemma}

策略是把 $U_{J,t}$ 上的复 MA 方程取极限。这是可行的，道德上因为复 MA 算子在势的 $C^0$-收敛下是弱连续的。在我们的情形有一个额外的微妙之处：这一列势定义在不同的流形上；对通常论证的修改在 [52, section 5.1] 中完成。

% [Tutor Note: 「补全循环」指：(17) 的实 MA 解 $\bar\phi_{J,0}$ ← 由极限凸势得到 ← 由复 MA 方程取极限；引理 6.9 补上「极限凸势确实解实 MA 方程」这一环，从而把 $Sk(X)$ 上的实 MA 方程与 $X_t$ 上的复几何真正接起来。$\Delta_J$ 内部（不是整个 $Sk(X)$）是关键限制。]

\subsubsection{与 NA 几何的关系}
先验极限策略并不显式地借助 NA 几何。其剩下的主要困难是延拓问题。基于 Fermat 例子的经验，我们预期延拓到环境环面簇上的环面度量是一项有用的技术，而该问题可能具有实质上的组合学面向。正如我们在第 5.5.6 节强调的，显式的凸势类对于直接表述实 MA 方程也不可或缺；而要处理更精细的问题------例如 $Sk(X)$ 上实 MA 度量的仿射结构与奇异集------很可能正需要这种表述（参看第 5.3.3 节的 Kontsevich--Soibelman 猜想）。

相反，NA 多复变势理论围绕 $X_K^{an}$ 上 NA 半正度量这一核心概念建立；经 Fubini--Study 逼近，这些度量可以 $C^0$-小误差地延拓为 $X_t$ 上的 Kähler 势。从这个意义上说，NA 多复变势理论可以看作延拓问题的一种伪装了的解答。然而，要与微分几何的应用接上，

% [原文 p.43]
还需要某种附加假设，例如 NA MA 与实 MA 比较性质。比较两种策略中的困难，我们猜测：证明 NA MA 与实 MA 比较性质需要对 NA 半正度量作更具体的刻画，或许具有显式的组合学特性。

% [Tutor Note: 本章正文至此结束（p43 顶部三行），紧接着是原文的 "References" 标题；原文既无末尾致谢（Acknowledgement 位于 p2，属 b12/tr-07 的页段），也无结论节，故本 chunk 不加 \section*{致谢}，也不写参考文献。全章最后一句把「证明比较性质」归结为「更具体地刻画 NA 半正度量」，与 5.6 节的问题 6、6.2.2 的延拓问题前后呼应。]
